\scanpage{69}
% archive foliation: 71
des syst. admissibles de disques \inserted{conformes} autour
des $a_i$, vers l'ens des syst.
de disques \deleted{conformes}\inserted{\unsure{standards|euclidiens}} dans les $T_{X,a_i}$,
est bijective. \deleted{Au niveau des
esp.}

Au niveau des espaces
modulaires de type Teichmüller,
cela donnerait donc une \deleted{\unsure{à}} structure
de fibré principal de groupe
\deleted{\diagram{\begin{tikzcd}[ampersand replacement=\&]TB_{g,\unsure{I|\nu}}\arrow[r,"\sim"]\&T_{g\nu}\end{tikzcd}}}
structural $(\mathbb R^{*+})^I$ : $TB_{g,I}$
sur $T_{g,I}$ $(\operatorname{card}I=\nu)$, \deleted{\unsure{--}}

Si \unsure{la relation} qui précédait
\deleted{(en passant aux log}
\deleted{longueurs} était vraie, cela signifierait
qu'on aurait une section canonique
de ce fibré, ce qui implique
aussi que pour un pt de $T\deleted{B}_{g,I}$,
i.e. une surface bd \inserted{ponctuée} ($\unsure{X_0}$-rigidifiée)
de type $(g,\nu)$, il y aurait un syst.
d'ombres can. dans les $T_{X,a_i}$, i.e. des
